\chapter{PiOS pi-orbital active space (menu 48)}
\label{ch:menu48}
\menulabel{48}

\section{Purpose}

The active space of a conjugated pi-system, built automatically
(Sayfutyarova \& Hammes-Schiffer 2019).  For an aromatic ring or a
conjugated chain the chemically meaningful active space is the pi space
--- the occupied and virtual MOs built from the atoms' out-of-plane p
orbitals --- and PiOS constructs it from a single-reference wavefunction:
the pi plane from the atom positions (Eqs.~(1)--(3)), the electron count
from the atomic connectivity (Section~IIC), the projection spectra and the
selection (Eqs.~(11)--(15)), and the Fock semicanonicalisation
(Eq.~(16)).

\section{What you need}

An \code{orca\_2json} export of the RHF calculation carrying the
\code{S-Matrix}, the \code{H-Matrix} and the \code{FockMatrix} pair
(\code{\{"MOCoefficients": true, "1elIntegrals": ["H", "S"],
"FockMatrix": ["J", "K"]\}}), the 0-based indices of the pi-system atoms,
and the system's mkl for the write-back.  The manifest schema is in
Chapter~\ref{ch:formats}.

\section{Walkthrough}

The example builds the pi space of benzene from the RHF/cc-pVDZ export
(the six ring carbons as the pi system): the plane, the per-atom electron
count, the projection spectra and the selected orbitals with their parent
SCF indices:

\lstinputlisting[style=fbkcode, firstline=55, lastline=69,
  title={The menu-48 example (plane, electron count, CAS(6e, 6o))}]
  {examples/expected/m48_pios.txt}

\lstinputlisting[style=fbkcode, firstline=191, lastline=197,
  title={The selected pi orbitals (semicanonical energies, parent SCF
  weights; same run)}]
  {examples/expected/m48_pios.txt}

The written \file{benzene\_rhf.pios.fbk.mkl} carries the partition
inactive occupied $|$ pi occupied $|$ pi virtual $|$ inactive virtual.

\section{What you get}

The report \code{<manifest>.pios.fbk.md} (the plane normal, the centroid,
the maximum out-of-plane deviation, the per-atom sigma-bond counts and pi
contributions, the two projection spectra with the selection marks, the
selected orbitals' energies and parent weights, and the boundary notes)
plus the gbw-ready mkl: \code{orca\_2mkl <name> -gbw}, then \code{!moread}
+ \code{\%moinp} in a \code{\%casscf} run with \code{nel N\_pi,e},
\code{norb |M|}.

\section{Conventions, cross-checks and boundaries (measured)}

\begin{readbox}
\begin{itemize}
  \item \textbf{The projection spectra carry the source's own validity
    reading.}  On the benzene fixture the occupied side selects 0.7789 /
    0.7649 / 0.7649 against 0.000 for every other occupied orbital, and
    the virtual side 1.000 / 1.000 / 1.000 against 0.235: the selection
    gaps are clean.  The projector trace closes exactly on both sides
    (2.309 $+$ 3.691 $=$ 6, the dimension of the p$_z$ space) --- the
    internal consistency of the construction.
  \item \textbf{The route is the calculation's own basis, not the source's
    IAO.}  The oriented p orbitals live on each atom's valence p shell
    (the outermost p shell of the export); the source builds them in an
    auxiliary MINAO basis through IAOs, which needs the cross-basis
    overlap ORCA does not export --- the same substitution the AVAS menu
    carries.  The measured consequence: the occupied pi block's
    eigenvalues read 0.78 rather than 1.0 (the missing part is the d
    polarisation a pure-p projection cannot carry), and against the
    CASSCF-optimised active space the guess reads SVD eigenvalues 0.9999 /
    0.9999 / 0.9999 / 0.7573 / 0.7573 / 0.6558 --- the occupied side
    essentially exact, the virtual gap equal to the converged pi*
    orbitals' own out-of-plane character.  The source's own benzene
    benchmark (0.9708--1.0) belongs to its IAO / aug-cc-pVTZ setup.
  \item \textbf{The written set is checked and deterministic.}  Orthonormal
    to the export's own precision in its overlap metric; the Fock matrix
    is assembled as H\,+\,J\,+\,K from the export's blocks (the menu-21
    cross-check); within degenerate clusters (symmetry-equivalent
    orbitals) any orthonormal basis is equally energy-ordered, and the
    written basis is fixed by a Gram-Schmidt of the parent SCF orbitals in
    index order, so the files do not depend on the BLAS build.
  \item \textbf{Engine validation.}  A CASSCF(6,6) started from the
    written benzene space converges to the same solution as the aufbau
    start ($-230.793818903$ vs $-230.793818898$~Eh) in 13
    macro-iterations against 7: on this fixture the aufbau window already
    contains the pi space, and the source's large gains belong to
    condensed-phase settings where the aufbau window fails outright.
    ORCA mechanics measured: the coefficients-only export after
    \code{!NoIter moread} works, while a \code{FockMatrix} conf against
    that gbw crashed \code{orca\_2json} (the densities sidecar of the run
    is not the SCF's).
\end{itemize}
\end{readbox}

\section{Refusals and related sections}

\begin{refusalbox}
\begin{itemize}
  \item An export without the S-Matrix, without coefficients, without the
    Fock family, or with fractional occupations (a CASSCF export) is
    refused with the next step; so are fewer than three pi atoms, a
    repeated or out-of-range atom index, an element outside the
    contribution rule (C, N, P --- or an N/P with an untabulated sigma
    count), a contribution list of the wrong length, and an electron count
    that is not an even number fitting the atom count.
  \item A non-planar atom set (maximum out-of-plane deviation beyond
    0.5~\AA{} after the plane fit) and a collinear set (degenerate moments
    of inertia) are refused: the pi direction is not defined on them.
\end{itemize}
\end{refusalbox}

Related: \menu{14} (AVAS --- the same projection idea onto a chosen target
shell), \menu{25} (AEGISS, the entropy-plus-projection selection on the
benzene platform), \menu{47} (the AOP rotation --- the other gbw-ready
initial-guess route), and Chapter~\ref{ch:criteria} for the criterion
table.
